% claim.tex -- AI4Math claim of record (AI-generated, 2026-09-28; instantiated
% from harness/templates/claim/claim.tex per SOP 08b).
% Machine-readable theorem anchors are the "% LEAN:" lines; scripts/paper-lint.sh
% and the claim audit check them. All Lean listings are verbatim, byte-identical
% contiguous excerpts of frozen artifacts of tasks/20260928-comb07-diagnotch:
%   statements  : formalized/01-comb07-diagnotch-statements.lean (143 lines, sha256
%                 16744fd37b7d8eb67965641f4e2b6dd8e7e3c2a994ea484766b6e3afb5861aca)
%   proof bodies: final/01-comb07-diagnotch-proved.lean (433 lines, sha256
%                 51721e9e315db12a35ca62218d8d458f32443e97f666e4a2ab452d0056f7cb43)
% Line ranges used: statement file 39-53, 55-87, 89-94, 97-116, 119-131, 133-136,
% 139-142; final file 147-172, 179-190, 215-229, 320-334. The same blocks and line
% ranges were verified byte-exact in the paper package papers/grid3n-diag-2notch
% (audit/listing/) and are reused here unchanged; statement-file theorem listings
% stop at the ":= by" line because the frozen snapshot carries sorry placeholders
% there, and the paper-lint 0-sorry sentinel forbids the word inside a listing.
% Verification: audit/listing-verify.txt (audit/inspect-tex.py --verify).
\documentclass[11pt]{article}

\usepackage[T1]{fontenc}
\usepackage[utf8]{inputenc}
\usepackage{amsmath,amssymb,amsthm}
\usepackage{graphicx}
\usepackage{hyperref}
\usepackage{xurl}
\setlength{\emergencystretch}{3em}
% lstlean.tex — Lean style for the listings package.
% Lineage: Jeremy Avigad (2015), adapted from Assia Mahboubi's SSR style;
% inspired by Olivier Verdier's unixode.sty for the Unicode literate table.
% No CTAN package or upstream repo exists; papers carry this file in their
% source package (cap set arXiv:1907.01449, sphere eversion arXiv:2210.07746).
% AI4Math adaptation (AI-generated, 2026-09-26): sorry/admit/sorryAx elevated
% to a dedicated error tier, matching the pipeline's 0-sorry constitution.
\usepackage{listings}
\usepackage{xcolor}

\definecolor{keywordcolor}{rgb}{0.7, 0.1, 0.1}   % red keywords
\definecolor{tacticcolor}{rgb}{0.0, 0.1, 0.6}    % blue tactics
\definecolor{commentcolor}{rgb}{0.4, 0.4, 0.4}   % grey comments
\definecolor{symbolcolor}{rgb}{0.0, 0.1, 0.6}    % blue symbols
\definecolor{sortcolor}{rgb}{0.1, 0.5, 0.1}      % green sorts
\definecolor{errorcolor}{rgb}{0.6, 0, 0}         % dark red: sorry/admit
\definecolor{stringcolor}{rgb}{0.5, 0.1, 0.5}    % purple strings

\lstdefinelanguage{lean}{
  % Anything between $ becomes LaTeX math mode
  mathescape=false,
  % Comments may or not include Latex commands
  texcl=false,
  % keywords, list taken from lean-syntax.txt
  morekeywords=[1]{import, prelude, protected, private, noncomputable,
    definition, meta, renaming, hiding, parameter, parameters, begin,
    constant, constants, lemma, variable, variables, theory, print,
    theorem, example, open, section, namespace, universe, universes, end,
    modifier, modifiers, using, export, exposing, axiom, inductive,
    coinductive, with, structure, class, instance, attribute, local, where,
    by, have, show, let, in, if, then, else, match, do, for, fun, assume,
    from, take, calc, include, omit, infix, infixl, infixr, notation,
    macro, syntax, elab, set_option, initialize, deriving, opaque, abbrev,
    mutual, partial, scoped, termination_by, decreasing_by},
  % Sorts
  morekeywords=[2]{Sort, Type, Prop},
  % tactics, list taken from lean-syntax.txt (tactic keywords)
  morekeywords=[3]{intro, intros, apply, exact, refine, rfl, simp, simp_all,
    simpa, simp_rw, rw, rewrite, erewrite, ring, ring_nf, omega, decide,
    norm_num, norm_cast, induction, cases, rcases, obtain, constructor,
    ext, funext, conv, congr, field_simp, linarith, nlinarith, positivity,
    tauto, trivial, aesop, use, by_cases, by_contra, push_neg, contrapose,
    symm, trans, assumption, contradiction, exfalso, specialize,
    generalize, revert, clear, rename_i, subst, unfold, delta, change,
    convert, split_ifs, interval_cases, fin_cases, zify, gcongr,
    nth_rewrite, suffices, generalizing, at, only, native_decide},
  % warnings - constitution tier: must never survive into a final proof
  morekeywords=[4]{sorry, admit, sorryAx},
  literate=
    {α}{{\ensuremath{\alpha}}}1
    {β}{{\ensuremath{\beta}}}1
    {γ}{{\ensuremath{\gamma}}}1
    {δ}{{\ensuremath{\delta}}}1
    {ε}{{\ensuremath{\varepsilon}}}1
    {ζ}{{\ensuremath{\zeta}}}1
    {η}{{\ensuremath{\eta}}}1
    {θ}{{\ensuremath{\theta}}}1
    {ι}{{\ensuremath{\iota}}}1
    {κ}{{\ensuremath{\kappa}}}1
    {λ}{{\ensuremath{\lambda}}}1
    {μ}{{\ensuremath{\mu}}}1
    {ν}{{\ensuremath{\nu}}}1
    {ξ}{{\ensuremath{\xi}}}1
    {π}{{\ensuremath{\pi}}}1
    {ρ}{{\ensuremath{\rho}}}1
    {σ}{{\ensuremath{\sigma}}}1
    {τ}{{\ensuremath{\tau}}}1
    {υ}{{\ensuremath{\upsilon}}}1
    {φ}{{\ensuremath{\varphi}}}1
    {χ}{{\ensuremath{\chi}}}1
    {ψ}{{\ensuremath{\psi}}}1
    {ω}{{\ensuremath{\omega}}}1
    {Γ}{{\ensuremath{\Gamma}}}1
    {Δ}{{\ensuremath{\Delta}}}1
    {Θ}{{\ensuremath{\Theta}}}1
    {Λ}{{\ensuremath{\Lambda}}}1
    {Ξ}{{\ensuremath{\Xi}}}1
    {Π}{{\ensuremath{\Pi}}}1
    {Σ}{{\ensuremath{\Sigma}}}1
    {Φ}{{\ensuremath{\Phi}}}1
    {Ψ}{{\ensuremath{\Psi}}}1
    {Ω}{{\ensuremath{\Omega}}}1
    {ℵ}{{\ensuremath{\aleph}}}1
    {∀}{{\ensuremath{\forall}}}1
    {∃!}{{\ensuremath{\exists}!}}1
    {∃}{{\ensuremath{\exists}}}1
    {∈}{{\ensuremath{\in}}}1
    {∉}{{\ensuremath{\notin}}}1
    {∘}{{\ensuremath{\circ}}}1
    {∩}{{\ensuremath{\cap}}}1
    {∪}{{\ensuremath{\cup}}}1
    {⊆}{{\ensuremath{\subseteq}}}1
    {⊂}{{\ensuremath{\subset}}}1
    {⊇}{{\ensuremath{\supseteq}}}1
    {⊃}{{\ensuremath{\supset}}}1
    {⊄}{{\ensuremath{\nsubseteq}}}1
    {∅}{{\ensuremath{\emptyset}}}1
    {∧}{{\ensuremath{\wedge}}}1
    {∨}{{\ensuremath{\vee}}}1
    {¬}{{\ensuremath{\neg}}}1
    {⊤}{{\ensuremath{\top}}}1
    {⊥}{{\ensuremath{\bot}}}1
    {↔}{{\ensuremath{\leftrightarrow}}}1
    {→}{{\ensuremath{\rightarrow}}}1
    {←}{{\ensuremath{\leftarrow}}}1
    {↑}{{\ensuremath{\uparrow}}}1
    {⇒}{{\ensuremath{\Rightarrow}}}1
    {⇐}{{\ensuremath{\Leftarrow}}}1
    {⟹}{{\ensuremath{\Longrightarrow}}}1
    {↦}{{\ensuremath{\mapsto}}}1
    {≤}{{\ensuremath{\leq}}}1
    {≥}{{\ensuremath{\geq}}}1
    {≠}{{\ensuremath{\neq}}}1
    {≈}{{\ensuremath{\approx}}}1
    {≡}{{\ensuremath{\equiv}}}1
    {≃}{{\ensuremath{\simeq}}}1
    {≅}{{\ensuremath{\cong}}}1
    {×}{{\ensuremath{\times}}}1
    {·}{{\ensuremath{\cdot}}}1
    {±}{{\ensuremath{\pm}}}1
    {⊕}{{\ensuremath{\oplus}}}1
    {⊗}{{\ensuremath{\otimes}}}1
    {⋆}{{\ensuremath{\star}}}1
    {▸}{{\ensuremath{\blacktriangleright}}}1
    {⟨}{{\ensuremath{\langle}}}1
    {⟩}{{\ensuremath{\rangle}}}1
    {⦃}{{\ensuremath{\{\!\mid}}}1
    {⦄}{{\ensuremath{\mid\!\}}}}1
    {⟦}{{\ensuremath{[\![}}}1
    {⟧}{{\ensuremath{]\!]}}}1
    {⌊}{{\ensuremath{\lfloor}}}1
    {⌋}{{\ensuremath{\rfloor}}}1
    {⌈}{{\ensuremath{\lceil}}}1
    {⌉}{{\ensuremath{\rceil}}}1
    {√}{{\ensuremath{\surd}}}1
    {∣}{{\ensuremath{\mid}}}1
    {∤}{{\ensuremath{\nmid}}}1
    {∎}{{\ensuremath{\blacksquare}}}1
    {⋯}{{\ensuremath{\cdots}}}1
    {…}{{\ensuremath{\ldots}}}1
    {∑}{{\ensuremath{\sum}}}1
    {∏}{{\ensuremath{\prod}}}1
    {⁻¹}{{\ensuremath{^{-1}}}}1
    {¹}{{\ensuremath{^1}}}1
    {²}{{\ensuremath{^2}}}1
    {³}{{\ensuremath{^3}}}1
    {ⁿ}{{\ensuremath{^n}}}1
    {ᵐ}{{\ensuremath{^m}}}1
    {₀}{{\ensuremath{_0}}}1
    {₁}{{\ensuremath{_1}}}1
    {₂}{{\ensuremath{_2}}}1
    {₃}{{\ensuremath{_3}}}1
    {₄}{{\ensuremath{_4}}}1
    {₅}{{\ensuremath{_5}}}1
    {ₙ}{{\ensuremath{_n}}}1
    {ₘ}{{\ensuremath{_m}}}1
    {ₖ}{{\ensuremath{_k}}}1
    {ᵢ}{{\ensuremath{_i}}}1
    {ⱼ}{{\ensuremath{_j}}}1
    {ℤ}{{\ensuremath{\mathbb{Z}}}}1
    {ℕ}{{\ensuremath{\mathbb{N}}}}1
    {ℚ}{{\ensuremath{\mathbb{Q}}}}1
    {ℝ}{{\ensuremath{\mathbb{R}}}}1
    {ℂ}{{\ensuremath{\mathbb{C}}}}1
    {𝔽}{{\ensuremath{\mathbb{F}}}}1
    {𝒫}{{\ensuremath{\mathcal{P}}}}1
    {∗}{{\ensuremath{\ast}}}1
    {•}{{\ensuremath{\bullet}}}1,
  % Comments
  morecomment=[l]{--},
  morecomment=[s]{/-}{-/},
  morestring=[b]",
  stringstyle=\color{stringcolor},
  keepspaces=true,
  keywordstyle=[1]\color{keywordcolor}\bfseries,
  keywordstyle=[2]\color{sortcolor}\bfseries,
  keywordstyle=[3]\color{tacticcolor}\bfseries,
  keywordstyle=[4]\color{errorcolor}\bfseries\underline,
  escapeinside={(*@}{@*)},
  columns=fullflexible,
  basicstyle=\ttfamily\small,
  commentstyle=\color{commentcolor}\itshape,
  breaklines=true,
  showstringspaces=false,
  frame=single,
  framesep=2mm,
  xleftmargin=4mm,
  numbers=left,
  numberstyle=\tiny\color{commentcolor},
  numbersep=6pt,
}

% Inline Lean code in prose: \lean{Int.ModEq.pow}
\newcommand{\lean}[1]{\lstinline[language=lean,basicstyle=\ttfamily\normalsize]|#1|}

\lstset{language=lean}

% --- local rendering patch (2026-09-28) ---
% tectonic/XeTeX does not apply the listings literate table to multi-byte
% characters, so the Unicode set used by the listings of THIS note is mapped to
% math glyphs as active characters via the standard \lccode/\lowercase idiom
% (ASCII-only source, no extra packages). Restricted to symbols outside xeCJK's
% own punctuation table: 00B7 also occurs in the frozen sources but is special
% punctuation that xeCJK sets up at load time -- making it active first aborts
% the run with "Control sequence ... already defined" (observed 2026-09-27); it
% renders through xeCJK + SimSun instead, as do the CJK comment characters and
% the fullwidth punctuation. Rendering only: listing source bytes are untouched,
% so the byte-exactness checks are unaffected, and the \ifdefined guard leaves
% the pdflatex path identical to the template.
\ifdefined\XeTeXversion
  {\lccode`\~="2022 \lowercase{\global\catcode`~=\active \global\def~{\ensuremath{\bullet}}}}
  {\lccode`\~="2115 \lowercase{\global\catcode`~=\active \global\def~{\ensuremath{\mathbb{N}}}}}
  {\lccode`\~="2124 \lowercase{\global\catcode`~=\active \global\def~{\ensuremath{\mathbb{Z}}}}}
  {\lccode`\~="2190 \lowercase{\global\catcode`~=\active \global\def~{\ensuremath{\leftarrow}}}}
  {\lccode`\~="2192 \lowercase{\global\catcode`~=\active \global\def~{\ensuremath{\rightarrow}}}}
  {\lccode`\~="2194 \lowercase{\global\catcode`~=\active \global\def~{\ensuremath{\leftrightarrow}}}}
  {\lccode`\~="2200 \lowercase{\global\catcode`~=\active \global\def~{\ensuremath{\forall}}}}
  {\lccode`\~="2203 \lowercase{\global\catcode`~=\active \global\def~{\ensuremath{\exists}}}}
  {\lccode`\~="2227 \lowercase{\global\catcode`~=\active \global\def~{\ensuremath{\wedge}}}}
  {\lccode`\~="2228 \lowercase{\global\catcode`~=\active \global\def~{\ensuremath{\vee}}}}
  {\lccode`\~="2261 \lowercase{\global\catcode`~=\active \global\def~{\ensuremath{\equiv}}}}
  {\lccode`\~="2264 \lowercase{\global\catcode`~=\active \global\def~{\ensuremath{\leq}}}}
  {\lccode`\~="2265 \lowercase{\global\catcode`~=\active \global\def~{\ensuremath{\geq}}}}
  {\lccode`\~="22A2 \lowercase{\global\catcode`~=\active \global\def~{\ensuremath{\vdash}}}}
  {\lccode`\~="27E8 \lowercase{\global\catcode`~=\active \global\def~{\ensuremath{\langle}}}}
  {\lccode`\~="27E9 \lowercase{\global\catcode`~=\active \global\def~{\ensuremath{\rangle}}}}
  % code points present in the comb-07 sources, missing from the template patch:
  % 21A6 (the subfamily/pattern-table arrows in the docstrings), 27FA (the "iff"
  % glyph in the T3 docstrings) and 2212 (the minus sign).
  {\lccode`\~="21A6 \lowercase{\global\catcode`~=\active \global\def~{\ensuremath{\mapsto}}}}
  {\lccode`\~="27FA \lowercase{\global\catcode`~=\active \global\def~{\ensuremath{\Longleftrightarrow}}}}
  {\lccode`\~="2212 \lowercase{\global\catcode`~=\active \global\def~{\ensuremath{-}}}}
\fi
\ifdefined\XeTeXversion
  \usepackage{xeCJK}
  \setCJKmainfont{SimSun}
\fi

\newtheorem{theorem}{Theorem}
\theoremstyle{definition}
\newtheorem{conjecture}{Conjecture}

\title{Tiling-count sequence of the $3\times n$ rectangle with two opposite
corner cells removed: a new recurrence, its interleave identity and parity laws
(Lean 4 machine-checked claim of record)}
\author{Aurora0134\thanks{Contact: via the GitHub account Aurora0134; ORCID to
be added at Zenodo upload.}}
\date{September 28, 2026}

\begin{document}
\maketitle

\begin{abstract}
Let $a(n)$ be the number of matchings (monomer--dimer coverings with
$1\times 2$ dominoes, the empty one included) of the $3\times n$ rectangle
with the top-right and the bottom-left cell removed:
$1,5,33,204,1266,7873,48882,\dots$. We state and machine-check, in Lean~4
(v4.34.0, mathlib v4.34.0), four theorems about the integer sequences defined
by the initial values and the recurrences fitted to this count: the main
order-$7$ recurrence and the interleaving of the two same-coefficient order-$7$
subfamily recurrences characterise each other, and the parities follow explicit
patterns (period $12$ modulo $2$ for the main sequence, period $6$ for the two
subsequences). Every proof compiles with no \texttt{sorry} and
\texttt{\#print axioms} reports only \{propext, Quot.sound,
Classical.choice\}. Novelty adjudication: no occupying record found (three
consistent rounds), value tier \emph{new sequence / new recurrence}, not new
mathematics; the perfect-matching sub-sequence at even indices is already
registered (OEIS A061278) and is not claimed. The combinatorial identification
is not kernel-proved and is stated as a conjecture. This note is a priority
claim of record, not a full paper.
\end{abstract}

\section{Introduction}

Four theorems about three integer sequences defined by initial values and
linear recurrences are reported here; the sequences come from counting
matchings of the $3\times n$ rectangle with the top-right and bottom-left cells
removed, and the recurrences were obtained by exact rational fitting of
enumerated data. The proofs were produced by the AI4Math pipeline, in which LLM
workers generate statements and proof scripts and the Lean 4 kernel is the sole
adjudicator at every gate: statements are frozen by hash before proving, and a
separate audit re-computes the symmetric difference, scans the whole file for
\texttt{sorry}/\texttt{admit} including comments, and re-prints the axiom
dependencies. Verification strength is stated below verbatim. This note is a
priority claim of record deposited so that a dated, citable public record
exists ahead of any narrative paper; the fuller paper for the same four
theorems is prepared separately and will reference this deposit. No literature
review is attempted here.

\section{Statement}\label{sec:statement}

Throughout, indexing is $0$-based in the Lean layer and $1$-based in the
counting layer: $\mathrm{adiag}\,n=a(n+1)$,
$\mathrm{onotch}\,k=a(2k+1)$ (odd positions) and
$\mathrm{enotch}\,k=a(2k+2)$ (even positions). The three sequences are
\emph{defined} by the displayed initial values and recurrences; the statement
layer contains no graphs, no matchings and no bijections (the type is
$\mathbb{N}\to\mathbb{Z}$ throughout).

\begin{lstlisting}[caption={Sequence \lean{adiag} (main sequence), verbatim from
the frozen statement file \texttt{formalized/01-comb07-diagnotch-statements.lean}
(lines 39--53, sha256 \texttt{16744fd3}).},label={lst:adiag}]
/-- **序列 adiag（主列，七阶递推）**。
初值 a(0..6) = 1, 5, 33, 204, 1266, 7873, 48882（0-based；对应 1-based a(1..7)）；
递推 a(n+7) = 4·a(n+6) + 15·a(n+5) − 5·a(n+4) − 19·a(n+3)
            + 9·a(n+2) + 2·a(n+1) − a(n)（n ≥ 0）。 -/
def adiag : ℕ → ℤ
  | 0 => 1
  | 1 => 5
  | 2 => 33
  | 3 => 204
  | 4 => 1266
  | 5 => 7873
  | 6 => 48882
  | n + 7 =>
      4 * adiag (n + 6) + 15 * adiag (n + 5) - 5 * adiag (n + 4)
        - 19 * adiag (n + 3) + 9 * adiag (n + 2) + 2 * adiag (n + 1) - adiag n
\end{lstlisting}

\begin{lstlisting}[caption={Sequences \lean{onotch} and \lean{enotch} (the two
subsequences, same coefficients), verbatim from the frozen statement file
(lines 55--87, sha256 \texttt{16744fd3}).},label={lst:sub}]
/-- **序列 onotch（1-based 奇位子列，与 enotch 同系数的七阶递推）**。
初值 o(0..6) = 1, 33, 1266, 48882, 1886600, 72804880, 2809532937
（= adiag 在下标 0,2,4,6,8,10,12 处之值）；
递推 o(k+7) = 46·o(k+6) − 303·o(k+5) + 671·o(k+4) − 519·o(k+3)
            + 167·o(k+2) − 22·o(k+1) + o(k)（k ≥ 0）。 -/
def onotch : ℕ → ℤ
  | 0 => 1
  | 1 => 33
  | 2 => 1266
  | 3 => 48882
  | 4 => 1886600
  | 5 => 72804880
  | 6 => 2809532937
  | k + 7 =>
      46 * onotch (k + 6) - 303 * onotch (k + 5) + 671 * onotch (k + 4)
        - 519 * onotch (k + 3) + 167 * onotch (k + 2) - 22 * onotch (k + 1) + onotch k

/-- **序列 enotch（1-based 偶位子列，与 onotch 同系数的七阶递推）**。
初值 e(0..6) = 5, 204, 7873, 303723, 11720017, 452270456, 17453021829
（= adiag 在下标 1,3,5,7,9,11,13 处之值）；
递推 e(k+7) = 46·e(k+6) − 303·e(k+5) + 671·e(k+4) − 519·e(k+3)
            + 167·e(k+2) − 22·e(k+1) + e(k)（k ≥ 0）。 -/
def enotch : ℕ → ℤ
  | 0 => 5
  | 1 => 204
  | 2 => 7873
  | 3 => 303723
  | 4 => 11720017
  | 5 => 452270456
  | 6 => 17453021829
  | k + 7 =>
      46 * enotch (k + 6) - 303 * enotch (k + 5) + 671 * enotch (k + 4)
        - 519 * enotch (k + 3) + 167 * enotch (k + 2) - 22 * enotch (k + 1) + enotch k
\end{lstlisting}

% LEAN: tasks/20260928-comb07-diagnotch :: adiag_interleave grade=完全证明
\begin{theorem}[interleaving identity]\label{thm:interleave}
Let \lean{adiag}, \lean{onotch}, \lean{enotch} be defined by the initial values
and the recurrences of Listings~\ref{lst:adiag} and~\ref{lst:sub}. Then for
every $k\in\mathbb{N}$,
\[
  \mathrm{adiag}\,(2k)=\mathrm{onotch}\,k
  \qquad\text{and}\qquad
  \mathrm{adiag}\,(2k+1)=\mathrm{enotch}\,k .
\]
Equivalently: the full order-$7$ recurrence and the interleaving of the two
same-coefficient order-$7$ subfamily recurrences characterise each other.
\end{theorem}

\begin{lstlisting}[caption={Theorem~\ref{thm:interleave}, verbatim from the
frozen statement file (lines 89--94, sha256 \texttt{16744fd3}); the listing
stops at the theorem's \texttt{:= by} line.},label={lst:t1stmt}]
/-- **T1（主攻·交织恒等式）**：对一切 k，adiag (2k) = onotch k 且
adiag (2k+1) = enotch k。即：主列的完整七阶递推 = 两条同系数
（46,−303,671,−519,167,−22,1）七阶子列递推（onotch/enotch）的交织，
两者互相刻画（0-based；对应 1-based 的 a(2k+1)=o(k)、a(2k+2)=e(k)）。 -/
theorem adiag_interleave (k : ℕ) :
    adiag (2 * k) = onotch k ∧ adiag (2 * k + 1) = enotch k := by
\end{lstlisting}

% LEAN: tasks/20260928-comb07-diagnotch :: adiag_mod2_period12 grade=完全证明
\begin{theorem}[main sequence, period $12$ modulo $2$]\label{thm:pat2}
For all $n,r\in\mathbb{N}$ with $n\bmod 12=r$,
\[
  \mathrm{adiag}\,n \bmod 2 = \mathrm{pat2}\,r,
\]
where \lean{pat2} is the twelve-entry pattern
$(1,1,1,0,0,1,0,1,0,1,0,0)$ of Listing~\ref{lst:pat2}. Equivalently,
$\mathrm{adiag}\,n$ is odd if and only if
$n\equiv 0,1,2,5,7,9\pmod{12}$.
\end{theorem}

\begin{lstlisting}[caption={Pattern table \lean{pat2} and
Theorem~\ref{thm:pat2}, verbatim from the frozen statement file (lines
97--116, sha256 \texttt{16744fd3}); the listing stops at the theorem's
\texttt{:= by} line.},label={lst:pat2}]
/-- **余数样式 pat2**（12 留数类 ↦ mod 2 余数）：
0↦1, 1↦1, 2↦1, 3↦0, 4↦0, 5↦1, 6↦0, 7↦1, 8↦0, 9↦1, 10↦0, 11↦0。 -/
def pat2 : ℕ → ℤ
  | 0 => 1
  | 1 => 1
  | 2 => 1
  | 3 => 0
  | 4 => 0
  | 5 => 1
  | 6 => 0
  | 7 => 1
  | 8 => 0
  | 9 => 1
  | 10 => 0
  | _ => 0

/-- **T2（伴随·主列 mod 2 周期 12 显式样式）**：adiag n 的 mod 2 余数仅依赖
n mod 12，样式为 pat2（即 adiag n 为奇 ⟺ n ≡ 0,1,2,5,7,9 (mod 12)）。 -/
theorem adiag_mod2_period12 (n r : ℕ) (hr : n % 12 = r) :
    adiag n % 2 = pat2 r := by
\end{lstlisting}

% LEAN: tasks/20260928-comb07-diagnotch :: onotch_mod2_period6 grade=完全证明
\begin{theorem}[odd-position subsequence, period $6$ modulo $2$]\label{thm:pat6o}
For all $k,r\in\mathbb{N}$ with $k\bmod 6=r$,
\[
  \mathrm{onotch}\,k \bmod 2 = \mathrm{pat6o}\,r,
\]
where \lean{pat6o} is the six-entry pattern $(1,1,0,0,0,0)$ of
Listing~\ref{lst:pat6}. Equivalently, $\mathrm{onotch}\,k$ is odd if and only
if $k\equiv 0$ or $1\pmod 6$.
\end{theorem}

% LEAN: tasks/20260928-comb07-diagnotch :: enotch_mod2_period6 grade=完全证明
\begin{theorem}[even-position subsequence, period $6$ modulo $2$]\label{thm:pat6e}
For all $k,r\in\mathbb{N}$ with $k\bmod 6=r$,
\[
  \mathrm{enotch}\,k \bmod 2 = \mathrm{pat6e}\,r,
\]
where \lean{pat6e} maps $1$ and $5$ to $0$ and every other residue to $1$
(Listing~\ref{lst:pat6}). Equivalently, $\mathrm{enotch}\,k$ is even if and
only if $k\equiv 1$ or $5\pmod 6$.
\end{theorem}

\begin{lstlisting}[caption={Pattern tables \lean{pat6o} and \lean{pat6e},
verbatim from the frozen statement file (lines 119--131, sha256
\texttt{16744fd3}).},label={lst:pat6}]
/-- **余数样式 pat6o**（onotch 侧 6 留数类 ↦ mod 2 余数）：
0↦1, 1↦1, 其余↦0（即 onotch k 为奇 ⟺ k ≡ 0 或 1 (mod 6)）。 -/
def pat6o : ℕ → ℤ
  | 0 => 1
  | 1 => 1
  | _ => 0

/-- **余数样式 pat6e**（enotch 侧 6 留数类 ↦ mod 2 余数）：
1↦0, 5↦0, 其余↦1（即 enotch k 为偶 ⟺ k ≡ 1 或 5 (mod 6)）。 -/
def pat6e : ℕ → ℤ
  | 1 => 0
  | 5 => 0
  | _ => 1
\end{lstlisting}

\begin{lstlisting}[caption={Theorem~\ref{thm:pat6o}, verbatim from the frozen
statement file (lines 133--136, sha256 \texttt{16744fd3}); the listing stops at
the theorem's \texttt{:= by} line.},label={lst:t3ostmt}]
/-- **T3 之 onotch（伴随·奇位子列 mod 2 周期 6 样式）**：onotch k 的 mod 2 余数
仅依赖 k mod 6，样式为 pat6o（即 onotch k 为奇 ⟺ k ≡ 0 或 1 (mod 6)）。 -/
theorem onotch_mod2_period6 (k r : ℕ) (hr : k % 6 = r) :
    onotch k % 2 = pat6o r := by
\end{lstlisting}

\begin{lstlisting}[caption={Theorem~\ref{thm:pat6e}, verbatim from the frozen
statement file (lines 139--142, sha256 \texttt{16744fd3}); the listing stops at
the theorem's \texttt{:= by} line.},label={lst:t3estmt}]
/-- **T3 之 enotch（伴随·偶位子列 mod 2 周期 6 样式）**：enotch k 的 mod 2 余数
仅依赖 k mod 6，样式为 pat6e（即 enotch k 为偶 ⟺ k ≡ 1 或 5 (mod 6)）。 -/
theorem enotch_mod2_period6 (k r : ℕ) (hr : k % 6 = r) :
    enotch k % 2 = pat6e r := by
\end{lstlisting}

\begin{conjecture}[combinatorial bridge, not proved here]\label{conj:bridge}
For every $n\ge 1$, the matching count $a(n)$ of the $3\times n$ rectangle with
the top-right and bottom-left cells removed equals
$\mathrm{adiag}\,(n-1)$; equivalently, $\mathrm{onotch}\,k=a(2k+1)$ and
$\mathrm{enotch}\,k=a(2k+2)$ for every $k$. This identification is not
kernel-proved, is not used by any theorem of this note, and rests on exact
enumeration by independently written programs with different coverage.
\end{conjecture}

\section{Proof}\label{sec:proof}

All four proofs are strong inductions on the index, with the seed window
discharged by computation (\texttt{decide}) and the step closed by congruence
algebra. Theorem~\ref{thm:interleave} first proves an auxiliary $14$-step
identity (\texttt{lifted}) by strong induction: for $n<7$ both sides are
explicit integers; for $n=m+7$ the definitional recurrence is unfolded at the
eight indices $m+7,\dots,m+21$, the seven induction hypotheses are rewritten
into the same $14$-step form, and a single \texttt{linear\_combination} of
these fifteen equations closes the step by cancellation. The main statement is
then a strong induction on $k$: for $k<7$ the fourteen equalities are concrete
integer identities; for $k=j+7$ the even half applies \texttt{lifted} at $2j$,
normalises the indices by \texttt{ring}, replaces the seven predecessors by the
induction hypotheses, and closes by the definitional unfolding of
\lean{onotch} (the odd half is identical, ending in \lean{enotch}).
Theorems~\ref{thm:pat2}--\ref{thm:pat6e} share one skeleton: strong induction,
one definitional unfolding whose coefficients collapse modulo $2$ to
$0,1,1,1,1,0,1$, and a finite check per residue class
(\texttt{mod\_cases} with \texttt{omega}/\texttt{decide}).

The complete proof file has $433$ lines and therefore exceeds one page; the
four contiguous verbatim excerpts below carry the distinctive argument shapes,
and the complete file ships in this deposit under \texttt{proofs/}.

\begin{lstlisting}[caption={Inductive step of the auxiliary lemma
\texttt{lifted}: the last unfolding \texttt{p14}, the seven induction
hypotheses \texttt{h0}--\texttt{h6} and the closing
\texttt{linear\_combination}, verbatim excerpt from
\texttt{final/01-comb07-diagnotch-proved.lean} (lines 147--172, sha256
\texttt{51721e9e}).},label={lst:t1close}]
        have p14 : adiag (m + 21) = 4 * adiag (m + 20) + 15 * adiag (m + 19)
            - 5 * adiag (m + 18) - 19 * adiag (m + 17) + 9 * adiag (m + 16)
            + 2 * adiag (m + 15) - adiag (m + 14) := by simp only [adiag]
        have h0 : adiag (m + 14) = 46 * adiag (m + 12) - 303 * adiag (m + 10)
            + 671 * adiag (m + 8) - 519 * adiag (m + 6) + 167 * adiag (m + 4)
            - 22 * adiag (m + 2) + adiag m := ihn m (by omega)
        have h1 : adiag (m + 15) = 46 * adiag (m + 13) - 303 * adiag (m + 11)
            + 671 * adiag (m + 9) - 519 * adiag (m + 7) + 167 * adiag (m + 5)
            - 22 * adiag (m + 3) + adiag (m + 1) := ihn (m + 1) (by omega)
        have h2 : adiag (m + 16) = 46 * adiag (m + 14) - 303 * adiag (m + 12)
            + 671 * adiag (m + 10) - 519 * adiag (m + 8) + 167 * adiag (m + 6)
            - 22 * adiag (m + 4) + adiag (m + 2) := ihn (m + 2) (by omega)
        have h3 : adiag (m + 17) = 46 * adiag (m + 15) - 303 * adiag (m + 13)
            + 671 * adiag (m + 11) - 519 * adiag (m + 9) + 167 * adiag (m + 7)
            - 22 * adiag (m + 5) + adiag (m + 3) := ihn (m + 3) (by omega)
        have h4 : adiag (m + 18) = 46 * adiag (m + 16) - 303 * adiag (m + 14)
            + 671 * adiag (m + 12) - 519 * adiag (m + 10) + 167 * adiag (m + 8)
            - 22 * adiag (m + 6) + adiag (m + 4) := ihn (m + 4) (by omega)
        have h5 : adiag (m + 19) = 46 * adiag (m + 17) - 303 * adiag (m + 15)
            + 671 * adiag (m + 13) - 519 * adiag (m + 11) + 167 * adiag (m + 9)
            - 22 * adiag (m + 7) + adiag (m + 5) := ihn (m + 5) (by omega)
        have h6 : adiag (m + 20) = 46 * adiag (m + 18) - 303 * adiag (m + 16)
            + 671 * adiag (m + 14) - 519 * adiag (m + 12) + 167 * adiag (m + 10)
            - 22 * adiag (m + 8) + adiag (m + 6) := ihn (m + 6) (by omega)
        linear_combination 4 * h6 + 15 * h5 - 5 * h4 - 19 * h3 + 9 * h2 + 2 * h1 - h0
          + p14 - 46 * p12 + 303 * p10 - 671 * p8 + 519 * p6 - 167 * p4 + 22 * p2 - p0
\end{lstlisting}

\begin{lstlisting}[caption={Even half of the inductive step of
Theorem~\ref{thm:interleave}, verbatim excerpt from
\texttt{final/01-comb07-diagnotch-proved.lean} (lines 179--190, sha256
\texttt{51721e9e}).},label={lst:t1main}]
      · rw [show 2 * (j + 7) = 2 * j + 14 by ring, lifted (2 * j)]
        rw [show 2 * j + 12 = 2 * (j + 6) by ring,
            show 2 * j + 10 = 2 * (j + 5) by ring,
            show 2 * j + 8 = 2 * (j + 4) by ring,
            show 2 * j + 6 = 2 * (j + 3) by ring,
            show 2 * j + 4 = 2 * (j + 2) by ring,
            show 2 * j + 2 = 2 * (j + 1) by ring]
        rw [(ih (j + 6) (by omega)).1, (ih (j + 5) (by omega)).1,
            (ih (j + 4) (by omega)).1, (ih (j + 3) (by omega)).1,
            (ih (j + 2) (by omega)).1, (ih (j + 1) (by omega)).1,
            (ih j (by omega)).1]
        simp only [onotch]
\end{lstlisting}

\begin{lstlisting}[caption={Core of Theorem~\ref{thm:pat2}: the modulo-$2$
collapse \texttt{h1} and the first residue branch, verbatim excerpt from
\texttt{final/01-comb07-diagnotch-proved.lean} (lines 215--229, sha256
\texttt{51721e9e}).},label={lst:t2core}]
        have h1 : adiag (j + 7) % 2
            = (adiag (j + 5) % 2 + adiag (j + 4) % 2 + adiag (j + 3) % 2
                + adiag (j + 2) % 2 + adiag j % 2) % 2 := by
          simp only [adiag]
          omega
        rw [h1, ih (j + 5) (by omega), ih (j + 4) (by omega),
            ih (j + 3) (by omega), ih (j + 2) (by omega), ih j (by omega)]
        mod_cases hj : j % 12
        · have hj0 : j % 12 = 0 := by simpa [Nat.ModEq] using hj
          have h5 : (j + 5) % 12 = 5 := by omega
          have h4 : (j + 4) % 12 = 4 := by omega
          have h3 : (j + 3) % 12 = 3 := by omega
          have h2 : (j + 2) % 12 = 2 := by omega
          have h7 : (j + 7) % 12 = 7 := by omega
          rw [hj0, h5, h4, h3, h2, h7]; decide
\end{lstlisting}

\begin{lstlisting}[caption={Core of Theorem~\ref{thm:pat6o}, verbatim excerpt
from \texttt{final/01-comb07-diagnotch-proved.lean} (lines 320--334, sha256
\texttt{51721e9e}); Theorem~\ref{thm:pat6e} has the same shape at lines
383--397.},label={lst:t3core}]
        have h1 : onotch (j + 7) % 2
            = (onotch (j + 5) % 2 + onotch (j + 4) % 2 + onotch (j + 3) % 2
                + onotch (j + 2) % 2 + onotch j % 2) % 2 := by
          simp only [onotch]
          omega
        rw [h1, ih (j + 5) (by omega), ih (j + 4) (by omega),
            ih (j + 3) (by omega), ih (j + 2) (by omega), ih j (by omega)]
        mod_cases hj : j % 6
        · have hj0 : j % 6 = 0 := by simpa [Nat.ModEq] using hj
          have h5 : (j + 5) % 6 = 5 := by omega
          have h4 : (j + 4) % 6 = 4 := by omega
          have h3 : (j + 3) % 6 = 3 := by omega
          have h2 : (j + 2) % 6 = 2 := by omega
          have h7 : (j + 7) % 6 = 1 := by omega
          rw [hj0, h5, h4, h3, h2, h7]; decide
\end{lstlisting}

\section{Verification evidence}

\begin{itemize}
\item Toolchain: Lean \texttt{v4.34.0} \cite{lean4}
(\texttt{leanprover/lean4:v4.34.0}) and mathlib4 \texttt{v4.34.0}, revision
\texttt{5ed29652} \cite{mathlib}, pinned. Reproduction: with that pin,
\texttt{lake env lean} on
\path{proofs/01-comb07-diagnotch-proved.lean} compiles the file with exit
code $0$ (measured $37.0$\,s for the assembled final file in the strict pass).
\item Statement integrity: the frozen statement file
\path{formalized/01-comb07-diagnotch-statements.lean} ($143$ lines) has
sha256 \texttt{16744fd3\allowbreak{}7b7d8eb6\allowbreak{}7965641f\allowbreak{}4e2b6dd8\allowbreak{}7e3c2a99\allowbreak{}4ea48476\allowbreak{}6b6e3afb\allowbreak{}5861aca};
the final file \path{final/01-comb07-diagnotch-proved.lean} ($433$ lines) has
sha256 \texttt{51721e9e\allowbreak{}315db12a\allowbreak{}35ca6221\allowbreak{}8d8d458f\allowbreak{}32443e97\allowbreak{}f666e4a2\allowbreak{}ab452d00\allowbreak{}56f7cb43}.
The final file was compared with the frozen file by symmetric difference over
all ten declarations (six \texttt{def} and four \texttt{theorem}): zero content
difference, only a reordering (the six definitions were moved in front of the
theorems); each theorem header matches byte-for-byte up to the
\texttt{:= by} line.
\item Kernel check: compiles with $0$ \texttt{sorry} (whole-file scan including
comments, case-insensitive, also for \texttt{admit}/\texttt{sorryAx});
\texttt{\#print axioms} output, verbatim (kept literal; the smaller monospace
width is typesetting only):
\begin{lstlisting}[basicstyle=\ttfamily\footnotesize,frame=none]
'adiag_interleave' depends on axioms: [propext, Classical.choice, Quot.sound]
'adiag_mod2_period12' depends on axioms: [propext, Quot.sound]
'onotch_mod2_period6' depends on axioms: [propext, Quot.sound]
'enotch_mod2_period6' depends on axioms: [propext, Quot.sound]
\end{lstlisting}
The strict re-compile emits no error and no warning. The four axiom lines are
the audit department's independent probe output; the gate-batch sidecar files
\path{axioms-adiag_interleave.txt} and its three siblings ship in this
deposit under \texttt{audit/} with the same four lines each.
\item Grading by the audit department: all four theorems \emph{complete
proof}; lowest grade reached for the case = \emph{complete proof}; no
scaffolding route was used.
\item Audit chain (originals in this deposit under \texttt{audit/}): final
adjudication \path{final-comb07-diagnotch.txt} (four independent checks,
including the symmetric difference and the independent axiom probe), the four
axiom printouts \path{axioms-adiag_interleave.txt},
\path{axioms-adiag_mod2_period12.txt},
\path{axioms-onotch_mod2_period6.txt},
\path{axioms-enotch_mod2_period6.txt}, and the full C9 query log
\path{c9-record.md}. The semantic-fidelity check (roundtrip re-translation)
of the underlying task ran on the same model node as its formalisation; that
independence weakness is mitigated, not removed, by kernel adjudication and
manual review.
\end{itemize}

\section{Novelty evidence}\label{sec:novelty}

Adjudication copied from the pipeline's exclusivity review (C9 gate), which ran
three rounds before the problem was accepted, by workers who re-derived the
data independently rather than trusting any external report: verdict
\textbf{no occupying record found} (three consistent rounds), value tier
\textbf{new sequence / new recurrence}. The tier is not a claim of new
mathematics and not a first-discovery claim: the order-$7$ recurrence is
machine-verified to be of a different spectrum than the order-$6$ recurrence of
the undefected mother family A033506 (neither satisfies the other's recurrence;
the two characteristic polynomials share one quintic factor), and the counting
method itself (state-matrix recursion) is documented prior art that would
produce this sequence as a special case. Summary by channel; the full query
log, including every zero-hit query, is in this deposit at
\texttt{audit/c9-record.md}.

\begin{center}
\resizebox{\textwidth}{!}{%
\begin{tabular}{lll}
\hline
channel & queries & result \\
\hline
OEIS, sequence search \cite{oeis} & 50+ independent strings (rounds 1--3) &
zero hits (full string, prefixes, tails, $\times2$, characteristic-polynomial
and numerator bare strings, mod-$3$ subsequences); anchors self-verified \\
OEIS, per-family records & 2 & full-column reads of A033506/A033507
\cite{oeis-a033506,oeis-a033507}: references and links cover complete grids
only, no defect variant \\
OpenAlex & 7 & zero same-form (top hits: dimer physics / knowledge graphs /
materials noise) \\
zbMATH & 5 & zero same-form \\
arXiv & 5 & zero same-form (\texttt{total=0} on two alias groups) \\
Crossref & 3 & zero same-form \\
Oh 2019, nearest neighbour & full text &
\cite{oh2019}: Theorem~4 (fixed monomer set) is a method with a different
statistic; no instance of this sequence; reference chain (Kong 2006,
Tzeng--Wu 2003, Wu 2006) and $12$ citing works read, none enumerates
defective-grid all-matchings \\
Lean ecosystem & 2 + tree walk &
zero hits; mathlib has no matching-count theorem; compfiles zero; sequencelib
structurally absent; public Lean code search on nine large terms:
\texttt{total=0} \\
\hline
\end{tabular}%
}
\end{center}

Two boundaries are explicit. The perfect-matching sub-sequence at even indices
($1,5,20,76,285,1065,\dots$) is already registered as OEIS A061278
\cite{oeis-a061278} and is not claimed as new; the enumeration anchors A033507
and A030186 \cite{oeis-a033507,oeis-a030186} are likewise registry entries, not
claims. And the residual gaps, kept from the audit trail: the general web
search layer was unavailable in the review environment, so the literature
evidence rests on the four structured channels above; Lundow 1996/1998 and Read
1982 (the mother-family founding literature) were not read in full text;
Guichard 2008 is behind a paywall; and the Math Stack Exchange API was
rate-limited. If any of these channels later returns a hit on the
defective-grid matching counts, the tier above weakens accordingly, and this
deposit will be amended by a later version rather than by editing this record.

\section{Scope and limitations}

Grade and boundary, stated as the audit states them. (a) The theorems cover
the recurrence-defined sequences only; their equality with the matching counts
of the notched rectangle is not kernel-proved and is
Conjecture~\ref{conj:bridge}. Its empirical support is enumeration by exact
programs: the main counting program covers $1\le n\le 50$, cross-checked
against a delete-and-branch enumerator for $1\le n\le 14$ and against a
brute-force backtracker for $1\le n\le 9$, with the undefected $3\times n$,
$4\times n$ and $2\times n$ families reproduced as external anchors; the
enumerators themselves are not verified, and the scripts and their outputs
stay in the working repository under \texttt{tasks/20260928-comb07-diagnotch/}.
(b) The recurrences are computational findings: Berlekamp--Massey fits over
$\mathbb{Q}$ with minimal order exactly $7$, verified digit-for-digit against
enumeration, not derived from a transfer-matrix argument. (c) The theorems are
about parity only; the mod-$4$ pattern (period $12$), the mod-$3$ patterns
(periods $968$ and $484$) and a general periodicity statement were examined
during claim selection and deliberately left unproved in this batch, and
appear in no theorem here. (d) The perfect-matching sub-sequence at even
indices is already registered as OEIS A061278 and is excluded from the novelty
claim; the value tier ``new sequence / new recurrence'' carries only the
spectral-difference meaning recorded above. (e) The gate-four sign-off of the
underlying task was exercised by the author's instruction to run this claim
lane on 2026-09-28, while the task report file still carries its draft header
(pending manual sign-off at gate four); both facts are recorded here as they
stand, and the task products themselves were not edited. (f) This note carries
no literature review and is a deposit record, not a survey; the narrative
paper for the same result is separate.

\section*{Data availability}

The deposit package for this claim contains the claim note (LaTeX source and
PDF), the Lean sources (\texttt{proofs/}: frozen statement file, final proof
file, \texttt{lean-toolchain}), and the verification and novelty-search records
(\texttt{audit/}). Zenodo: the DOI is reserved at upload by the author and is
recorded in \texttt{metadata/README.md} of this package once created; at the
time of writing it is an unfilled field, \texttt{RESERVED-DOI}, not a value
(see the upload steps in \texttt{UPLOAD.md} in the source repository, which is
not shipped with this deposit). Public mirror of this deposit: the claim lane's repository
\url{https://github.com/Aurora0134/ai4math-claims}, directory
\texttt{grid3n-diag-2notch/} (pushed under the author's instruction of
2026-09-28; the commit and timestamp are recorded in the source repository's
\texttt{claims/grid3n-diag-2notch/card.md}).
Claim note under CC BY 4.0, code under MIT.

\section*{Use of generative AI}

(a) AI performed: candidate-problem screening and the three-round exclusivity
review, the numerical recomputation and claim selection of phase 0,
autoformalisation of the four statements and the roundtrip semantic check,
proof search and proof-script writing, the prose translation of the proofs,
and the assembly of this note including its verbatim listings. (b) Model, node
and budget: all task-stage sampling ran on a single node, wire-id
\texttt{anthropic/a6api-main/kimi-k3}, with no cross-node switching. One drift
is disclosed: one audit dispatch on 2026-09-28 (the F1 closure re-review)
carried a model-detection reading of
\texttt{anthropic/stepfun/step-5-preview}; the ledger's accounting node stayed
\texttt{anthropic/a6api-main/kimi-k3} and the drift is disclosed rather than
reconciled by substitution. The task consumed $11$ agent-run dispatches of a
permitted $20$ and $1$ \texttt{llm-call}
of a permitted $6$ (the roundtrip semantic judge; two earlier network failures
on the same node were retried and not counted), numbers copied verbatim from
the task ledger \texttt{tasks/20260928-comb07-diagnotch/budget.log} (a
source-repository file, not shipped with this package). This note consumed $0$
pipeline \texttt{llm-call} invocations and $1$ agent-run dispatch (this
assembly) against caps of $4$ and $4$ respectively, and its LaTeX compilation
used the rounds recorded in \texttt{audit/compile.txt}; the assembly session's
endpoint was \texttt{zcode}, whose model-detection reading of 2026-09-28
resolved to \texttt{stepfun/step-5-preview}, which is disclosed rather than
replaced by the task node's name. (c) Final adjudication: all statements and
proofs reported here were checked by the Lean 4 kernel: every proof compiles
with no \texttt{sorry}, and \texttt{\#print axioms} reports only \{propext,
Quot.sound, Classical.choice\} --- exactly this set for \lean{adiag_interleave},
and \{propext, Quot.sound\} for \lean{adiag_mod2_period12},
\lean{onotch_mod2_period6} and \lean{enotch_mod2_period6}. (d) The AI
system is not an author; the human author reviewed the evidence and is
responsible for all content.

\section*{Author contributions}

CRediT-style: the human author adjudicated topic selection, statement
semantics, the sign-off and the deposit decision; the AI4Math pipeline (an
LLM) generated the candidate propositions, the proof searches and the text
drafts. The AI system is not listed as an author.

\begin{thebibliography}{8}
% Reused from papers/grid3n-diag-2notch/paper.tex, where each entry was verified
% on 2026-09-28 before being typed in: Crossref REST metadata for the two DOIs
% (title and page range matched exactly), the arXiv export API for the preprint
% (title matched), and the OEIS JSON API for the four sequence entries (name and
% data matched; the A033506 signature link was read). Raw responses:
% papers/grid3n-diag-2notch/audit/bib-verify/ (source repository, not shipped
% with this deposit). No entry is quoted from memory; only entries actually
% cited above are listed.
\bibitem{lean4}
L.~de Moura and S.~Ullrich.
\newblock The Lean 4 theorem prover and programming language.
\newblock \emph{Automated Deduction -- CADE-28}, LNCS, pp.~625--635, 2021.
\url{https://doi.org/10.1007/978-3-030-79876-5_37}

\bibitem{mathlib}
The mathlib Community.
\newblock The Lean mathematical library.
\newblock \emph{CPP 2020: Proceedings of the 9th ACM SIGPLAN International
Conference on Certified Programs and Proofs}, pp.~367--381, 2020.
\url{https://doi.org/10.1145/3372885.3373824}

\bibitem{oh2019}
S.~Oh.
\newblock State matrix recursion method and monomer--dimer problem.
\newblock \emph{arXiv:1901.07847 [math.CO]}, 2019.
\url{https://arxiv.org/abs/1901.07847}

\bibitem{oeis-a033506}
The OEIS Foundation.
\newblock Entry A033506: number of matchings in graph $P_3 \times P_n$
(signature $(4,14,0,-10,0,1)$).
\newblock \emph{The On-Line Encyclopedia of Integer Sequences}.
\url{https://oeis.org/A033506}

\bibitem{oeis-a033507}
The OEIS Foundation.
\newblock Entry A033507: number of matchings in graph $P_4 \times P_n$.
\newblock \emph{The On-Line Encyclopedia of Integer Sequences}.
\url{https://oeis.org/A033507}

\bibitem{oeis-a030186}
The OEIS Foundation.
\newblock Entry A030186: number of matchings in graph $P_2 \times P_n$
($a(n)=3a(n-1)+a(n-2)-a(n-3)$).
\newblock \emph{The On-Line Encyclopedia of Integer Sequences}.
\url{https://oeis.org/A030186}

\bibitem{oeis-a061278}
The OEIS Foundation.
\newblock Entry A061278: $a(n)=5a(n-1)-5a(n-2)+a(n-3)$ --- the
perfect-matching sub-sequence at even $n$ of this note's graph family,
already registered and not claimed as new.
\newblock \emph{The On-Line Encyclopedia of Integer Sequences}.
\url{https://oeis.org/A061278}

\bibitem{oeis}
The OEIS Foundation.
\newblock \emph{The On-Line Encyclopedia of Integer Sequences}.
\newblock \url{https://oeis.org/} --- searched for the sequences, their
residue-class subsequences, their characteristic polynomials and their
generating-function numerators, with no entry found (see
Section~\ref{sec:novelty}).
\end{thebibliography}

\end{document}
